% GENERATED FILE. Edit canonical lessons, then run npm run generate:books.
% canonical-source-hash: fnv1a64:8cb02b0aa14954bf
% canonical-lessons: HI-C362-niskars, HI-C362-ausat, HI-C362-ubhaarna, HI-C362-vislesan, HI-C362-pariksan

\chapter{What the Figures Say}
\label{ch:hi-what-the-figures-say}
\hlchaptermodality{362}

\begin{chapteropening}
\textbf{By the end of this chapter:} \emph{I can say what a test showed.}

Complete the last lesson of chapter 362: I can say what a test showed.
\end{chapteropening}

\section[\texorpdfstring{\hi{निष्कर्ष} (niṣkarṣ)}{निष्कर्ष (niṣkarṣ)}]{\hi{निष्कर्ष} (niṣkarṣ) — a conclusion}

\label{lesson:HI-C362-niskars}



\begin{quote}
Before the new one: say the Hindi for to gauge, then the Hindi for a laboratory.
\end{quote}

\subsection*{You'll want to know: \hi{निष्कर्ष}}
\textbf{\hi{निष्कर्ष}} — \emph{niṣkarṣ} — \textquotedblleft{}a conclusion\textquotedblright{}.

\begin{quote}
\hi{बहस} \hi{के} \hi{बाद} \hi{हम} \hi{एक} \hi{निष्कर्ष} \hi{पर} \hi{पहुँचे।}
\end{quote}

\begin{quote}
\emph{bahas ke bād ham ek niṣkarṣ par pahuñce.}
\end{quote}

\begin{quote}
After the debate we reached a conclusion.
\end{quote}

Masculine.

\begin{cousinweb}
\textbf{\hi{निष्कर्ष}} (\emph{niṣkarṣ}) is a tatsama, borrowed from Sanskrit \emph{niṣkarṣa}, \textquotedblleft{}drawing out, extracting\textquotedblright{} (Monier-Williams). Hindi uses it for the conclusion drawn out of an argument.
\end{cousinweb}

\begin{grammarlens}[title={What you've built}]
One more word for giving reasons and opinions.
\end{grammarlens}

\subsection*{Your turn}
\begin{itemize}
\raggedright
  \item \emph{Say it:} \emph{niṣkarṣ}
  \item \emph{Say it:} \emph{niṣkarṣ}, once more
  \item \emph{Say it:} \emph{bahas ke bād ham ek niṣkarṣ par pahuñce.} — After the debate we reached a conclusion.
  \item \emph{Recall:} say the Hindi for to gauge, then the Hindi for a laboratory, then say \emph{niṣkarṣ} again
\end{itemize}

\begin{tcolorbox}[breakable,colback=teal!4,colframe=teal!35!black,title={Before you move on}]
What does \hi{निष्कर्ष} mean? (\textquotedblleft{}A conclusion\textquotedblright{}.) Say it once more.
\end{tcolorbox}
\section[\texorpdfstring{\hi{औसत} (ausat)}{औसत (ausat)}]{\hi{औसत} (ausat) — an average}

\label{lesson:HI-C362-ausat}



\begin{quote}
Before the new one: say the Hindi for a laboratory, then the Hindi for a conclusion.
\end{quote}

\subsection*{You'll want to know: \hi{औसत}}
\textbf{\hi{औसत}} — \emph{ausat} — \textquotedblleft{}an average\textquotedblright{}.

\begin{quote}
\hi{पिछले} \hi{साल} \hi{बारिश} \hi{औसत} \hi{से} \hi{कम} \hi{हुई।}
\end{quote}

\begin{quote}
\emph{pichle sāl bāriś ausat se kam huī.}
\end{quote}

\begin{quote}
Last year the rain was below average.
\end{quote}

Masculine.

\begin{cousinweb}
\textbf{\hi{औसत}} (\emph{ausat}) is usually traced to Arabic \emph{awsaṭ}, \textquotedblleft{}middlemost\textquotedblright{}, which reached Hindi through Persian.
\end{cousinweb}

\begin{grammarlens}[title={What you've built}]
One more word for giving reasons and opinions.
\end{grammarlens}

\subsection*{Your turn}
\begin{itemize}
\raggedright
  \item \emph{Say it:} \emph{ausat}
  \item \emph{Say it:} \emph{ausat}, once more
  \item \emph{Say it:} \emph{pichle sāl bāriś ausat se kam huī.} — Last year the rain was below average.
  \item \emph{Recall:} say the Hindi for a laboratory, then the Hindi for a conclusion, then say \emph{ausat} again
\end{itemize}

\begin{tcolorbox}[breakable,colback=teal!4,colframe=teal!35!black,title={Before you move on}]
What does \hi{औसत} mean? (\textquotedblleft{}An average\textquotedblright{}.) Say it once more.
\end{tcolorbox}
\section[\texorpdfstring{\hi{उभारना} (ubhārnā)}{उभारना (ubhārnā)}]{\hi{उभारना} (ubhārnā) — to bring out, to raise (a question)}

\label{lesson:HI-C362-ubhaarna}



\begin{quote}
Before the new one: say the Hindi for a conclusion, then the Hindi for an average.
\end{quote}

\subsection*{You'll want to know: \hi{उभारना}}
\textbf{\hi{उभारना}} — \emph{ubhārnā} — \textquotedblleft{}to bring out, to raise (a question)\textquotedblright{}.

\begin{quote}
\hi{इस} \hi{प्रयोग} \hi{ने} \hi{कई} \hi{नए} \hi{सवाल} \hi{उभारे।}
\end{quote}

\begin{quote}
\emph{is prayog ne kaī nae savāl ubhāre.}
\end{quote}

\begin{quote}
This experiment raised several new questions.
\end{quote}

Past, with \hi{ने} (he or she …): \hi{उसने} \hi{वही} \hi{सवाल} \hi{उभारा}, \emph{usne vahī savāl ubhārā}.

\begin{cousinweb}
\textbf{\hi{उभारना}} (\emph{ubhārnā}) is a tadbhava, inherited from Sanskrit \emph{udbharati}, \textquotedblleft{}raises up\textquotedblright{} (CDIAL 2038).
\end{cousinweb}

\begin{grammarlens}[title={What you've built}]
One more word for giving reasons and opinions.
\end{grammarlens}

\subsection*{Your turn}
\begin{itemize}
\raggedright
  \item \emph{Say it:} \emph{ubhārnā}
  \item \emph{Say it:} \emph{ubhārnā}, once more
  \item \emph{Say it:} \emph{is prayog ne kaī nae savāl ubhāre.} — This experiment raised several new questions.
  \item \emph{Recall:} say the Hindi for a conclusion, then the Hindi for an average, then say \emph{ubhārnā} again
\end{itemize}

\begin{tcolorbox}[breakable,colback=teal!4,colframe=teal!35!black,title={Before you move on}]
What does \hi{उभारना} mean? (\textquotedblleft{}To bring out, to raise (a question)\textquotedblright{}.) Say it once more.
\end{tcolorbox}
\section[\texorpdfstring{\hi{विश्लेषण} (viśleṣaṇ)}{विश्लेषण (viśleṣaṇ)}]{\hi{विश्लेषण} (viśleṣaṇ) — an analysis}

\label{lesson:HI-C362-vislesan}



\begin{quote}
Before the new one: say the Hindi for an average, then the Hindi for to bring out.
\end{quote}

\subsection*{You'll want to know: \hi{विश्लेषण}}
\textbf{\hi{विश्लेषण}} — \emph{viśleṣaṇ} — \textquotedblleft{}an analysis\textquotedblright{}.

\begin{quote}
\hi{शिक्षक} \hi{ने} \hi{कहानी} \hi{का} \hi{विश्लेषण} \hi{किया।}
\end{quote}

\begin{quote}
\emph{śikṣak ne kahānī kā viśleṣaṇ kiyā.}
\end{quote}

\begin{quote}
The teacher analysed the story.
\end{quote}

Masculine.

\begin{cousinweb}
\textbf{\hi{विश्लेषण}} (\emph{viśleṣaṇ}) is a tatsama, borrowed from Sanskrit \emph{viśleṣaṇa}, \textquotedblleft{}separation\textquotedblright{} (Monier-Williams). Hindi uses the idea of taking apart for an analysis.
\end{cousinweb}

\begin{grammarlens}[title={What you've built}]
One more word for giving reasons and opinions.
\end{grammarlens}

\subsection*{Your turn}
\begin{itemize}
\raggedright
  \item \emph{Say it:} \emph{viśleṣaṇ}
  \item \emph{Say it:} \emph{viśleṣaṇ}, once more
  \item \emph{Say it:} \emph{śikṣak ne kahānī kā viśleṣaṇ kiyā.} — The teacher analysed the story.
  \item \emph{Recall:} say the Hindi for an average, then the Hindi for to bring out, then say \emph{viśleṣaṇ} again
\end{itemize}

\begin{tcolorbox}[breakable,colback=teal!4,colframe=teal!35!black,title={Before you move on}]
What does \hi{विश्लेषण} mean? (\textquotedblleft{}An analysis\textquotedblright{}.) Say it once more.
\end{tcolorbox}
\section[\texorpdfstring{\hi{परीक्षण} (parīkṣaṇ)}{परीक्षण (parīkṣaṇ)}]{\hi{परीक्षण} (parīkṣaṇ) — a test, a trial (in research)}

\label{lesson:HI-C362-pariksan}



\begin{quote}
Before the new one: say the Hindi for to bring out, then the Hindi for an analysis.
\end{quote}

\subsection*{You'll want to know: \hi{परीक्षण}}
\textbf{\hi{परीक्षण}} — \emph{parīkṣaṇ} — \textquotedblleft{}a test, a trial (in research)\textquotedblright{}.

\begin{quote}
\hi{नई} \hi{दवा} \hi{का} \hi{परीक्षण} \hi{चल} \hi{रहा} \hi{है।}
\end{quote}

\begin{quote}
\emph{naī davā kā parīkṣaṇ cal rahā hai.}
\end{quote}

\begin{quote}
The new medicine is being tested.
\end{quote}

Masculine.

\begin{cousinweb}
\textbf{\hi{परीक्षण}} (\emph{parīkṣaṇ}) is a tatsama, borrowed from Sanskrit \emph{parīkṣaṇa}, \textquotedblleft{}trying, testing\textquotedblright{} (Monier-Williams).
\end{cousinweb}

\begin{grammarlens}[title={What you've built}]
One more word for giving reasons and opinions.
\end{grammarlens}

\subsection*{Your turn}
\begin{itemize}
\raggedright
  \item \emph{Say it:} \emph{parīkṣaṇ}
  \item \emph{Say it:} \emph{parīkṣaṇ}, once more
  \item \emph{Say it:} \emph{naī davā kā parīkṣaṇ cal rahā hai.} — The new medicine is being tested.
  \item \emph{Recall:} say the Hindi for to bring out, then the Hindi for an analysis, then say \emph{parīkṣaṇ} again
\end{itemize}

\begin{tcolorbox}[breakable,colback=teal!4,colframe=teal!35!black,title={Before you move on}]
What does \hi{परीक्षण} mean? (\textquotedblleft{}A test, a trial (in research)\textquotedblright{}.) Say it once more.
\end{tcolorbox}
